\section*{\texorpdfstring{\S\CurrentExerciseGroup}{Section \CurrentExerciseGroup}}
\phantomsection
\addcontentsline{toc}{section}{\texorpdfstring{Exercises for \S\CurrentExerciseGroup}{Exercises for Section \CurrentExerciseGroup}}

¶ 1) Let I be a semi-open interval \([a, b]\) in \(\mathbf{R}\) , and let \(F = R^I\) be the space of all mappings of I into \(\mathbf{R}\) , equipped with the topology of pointwise convergence; for every \(x \in I\) , denote by \(\mathbf{f}(x)\) the mapping \(t \mapsto |x - t|\) of I into \(\mathbf{R}\) , which is an element of \(F\) ; the mapping \(f\) of I into \(F\) is continuous. Show that \(f\) is left-differentiable at every point of I, but the left-derivative \(\mathbf{f}_l'\) is a function (with values in \(F\) ) that is not measurable for Lebesgue measure, even though it is the pointwise limit of a sequence of continuous functions and, for each \(t \in I\) , the function \(\text{pr}_t \circ \mathbf{f}_l'\) is a measurable numerical function (note that \(\mathbf{f}_l'\) is not right-continuous at any point of I, and make use of Exer. 1 of GT, IV, §2).

\begin{enumerate}
\def\labelenumi{\arabic{enumi})}
\setcounter{enumi}{1}
\tightlist
\item
  Let \(\nu\) be Lebesgue measure on \(\mathbf{R}\) , \(g\) a continuous mapping of \(\mathbf{R}\) into \([0,1]\) with support contained in \([-1,2[\) and equal to 1 on \([0,1]\) , and let \(\mu\) be the measure \(g \cdot \nu\) on \(\mathbf{R}\) . The set \(\mathrm{H}_0 \subset [0,1]\) defined in Exer. 8 of §4 is not \(\mu\) -measurable, but the set \(\mathrm{H} = \mathrm{H}_0 - 2\) is \(\mu\) -negligible.
\end{enumerate}

\begin{enumerate}
\def\labelenumi{\alph{enumi})}
\item
  If one sets \(f(x) = x - 2\) , then \(f\) is continuous and \(\varphi_{\mathrm{H}}\) is \(\mu\) -measurable, but the composed function \(\varphi_{\mathrm{H}} \circ f\) is not \(\mu\) -measurable.
\item
  If one sets \(h(x) = x + 2\) , the image \(h(\mathrm{H})\) of the \(\mu\) -measurable set \(\mathrm{H}\) under the continuous function \(h\) is not \(\mu\) -measurable.
\end{enumerate}

\begin{enumerate}
\def\labelenumi{\arabic{enumi})}
\setcounter{enumi}{2}
\tightlist
\item
  Let \(f\) be a measurable mapping of \(X\) into a topological space \(F\) , and \(g\) a lower semi-continuous numerical function on \(F\) ; show that \(g \circ f\) is measurable.
\end{enumerate}

¶ 4) Let \(\mu\) be Lebesgue measure on \(X = [0,1[\) , and let \((H_n)\) be a partition of \(X\) into an infinite sequence of sets having the power of the continuum none of which is measurable, and which are such that for every union \(H\) of a finite number of sets \(H_n\) , \(\mu_*(H) = 0\) (§4, Exer. 8). Let \(\sigma_n\) be a bijection of the interval \(]1/(n+1),1/n]\) onto \(H_n\) . For every number \(y\) such that \(0 < y \leqslant 1\) , let \(n\) be the integer such that \(1/(n+1) < y \leqslant 1/n\) . Define \(f_y\) to be the characteristic function of the set reduced to the point \(\sigma_n(y)\) ; for every \(x \in X\) , \(f_y(x)\) tends to 0 as \(y\) tends to 0. Show that there does not exist any compact set \(K \subset X\) of measure \textgreater0 such that \(f_y\) tends uniformly to 0 on \(K\) .

\begin{enumerate}
\def\labelenumi{\arabic{enumi})}
\setcounter{enumi}{4}
\tightlist
\item
  Let \((f_{mn})\) be a double sequence of measurable mappings of \(X\) into a metrizable space \(F\) . Assume that for every \(m\) , the sequence \((f_{mn})_{n\geqslant 1}\) converges locally almost everywhere to a function \(g_{m}\) , and that the sequence \((g_{m})\) converges locally almost everywhere to a function \(h\) . Show that for every compact subset \(K\) of \(X\) , there exist two strictly increasing sequences \((m_k),(n_k)\) of integers \(>0\) , such that the sequence of functions \(f_{m_k,n_k}\) converges to \(h\) almost everywhere in \(K\) . (Note that for every \(\varepsilon >0\) , there exists a compact set \(K_1\subset K\) such that \(\mu (K - K_1)\leqslant \varepsilon\) and such that in \(K_{1}\) the sequence \((g_{m})\) and each of the sequences \((f_{mn})_{n\geqslant 1}\) are uniformly convergent.)
\end{enumerate}

¶ 6) For every integer \(n \geqslant 1\) , let \(f_n(x) = [2^n x] - 2[2^{n-1}x]\) . In the compact space of mappings of \(\mathbf{R}\) into \(\{0,1\}\) (equipped with the topology of pointwise convergence), let \(f\) be a cluster point of the sequence \((f_n)\) . Show that, for every dyadic number \(r\) , \(f(r+x) = f(x)\) for all \(x \in \mathbf{R}\) and \(f(r-x) = 1 - f(x)\) for every \(x \in \mathbf{R}\) different from a dyadic number. From this, deduce that for the Lebesgue measure \(\mu\) , \(f\) is not measurable. (Argue by contradiction; let \(A\) be the set of \(x \in [0,1]\) such that \(f(x) = 1\) , and suppose that \(A\) is measurable and that \(\mu(A) = \alpha > 0\) ; show that there exists a set \(I\) , the finite union of open intervals contained in \(]0,1[\) , such that \(\mu(I \cap \mathbb{C}A) \leqslant \alpha/4\) and \(\mu(A \cap \mathbb{C}I) \leqslant \alpha/4\) ; on considering the intervals contained in \(I\) and of the form \(]k/2^n, (k+1)/2^n[\) , show that one obtains a contradiction with the relation \(f(r-x) = 1 - f(x)\) for \(r\) dyadic and \(x\) non-dyadic.)

\begin{enumerate}
\def\labelenumi{\arabic{enumi})}
\setcounter{enumi}{6}
\tightlist
\item
  Let \(\mathbf{X}\) be the interval \([0,1]\) in \(\mathbf{R}\) , and \(\mathbf{F}\) the Hilbert space having an orthonormal basis \((\mathbf{e}_t)_{0\leqslant t\leqslant 1}\) equipotent to \(\mathbf{X}\) .
\end{enumerate}

\begin{enumerate}
\def\labelenumi{\alph{enumi})}
\item
  Show that the mapping \(\mathbf{f}\) of \(X\) into \(F\) , such that \(\mathbf{f}(t) = \mathbf{e}_t\) for \(0 \leqslant t \leqslant 1\) , is not measurable for Lebesgue measure, but that the inverse image under \(\mathbf{f}\) of every closed ball in \(F\) is measurable and, for every continuous linear form \(\mathbf{a}'\) on \(F\) , \(\langle \mathbf{f}, \mathbf{a}' \rangle\) is negligible.
\item
  Let H be a non-measurable set in X (§4, Exer. 8); show that if \(\mathbf{g} = \mathbf{f}\varphi_{\mathrm{H}}\) then the function \(\langle \mathbf{g},\mathbf{a}'\rangle\) is negligible for every continuous linear form \(\mathbf{a}'\) on F, but the numerical function \(|\mathbf{g}|\) is not measurable.
\end{enumerate}

\begin{enumerate}
\def\labelenumi{\arabic{enumi})}
\setcounter{enumi}{7}
\item
  Let \(\mathbf{F}\) be a metrizable locally convex space, \(\mathbf{f}\) a mapping of \(\mathbf{X}\) into \(\mathbf{F}\) satisfying the conditions \(a)\) and \(b)\) of No.~5, Cor. 1 of Prop. 10; show that \(\mathbf{f}\) is measurable.
\item
  Let \(\mu\) be Lebesgue measure on \(X = [0,1]\) ; denote by \(F\) the vector space over \(\mathbf{R}\) of \(\mu\) -measurable finite numerical functions on \(X\) , equipped with the topology of pointwise convergence, which makes it a Hausdorff locally convex space.
\end{enumerate}

\begin{enumerate}
\def\labelenumi{\alph{enumi})}
\item
  Show that there exists in \(\mathbf{F}\) a countable dense subset (consider a dense sequence in the Banach space \(\mathcal{C}(\mathbf{X};\mathbf{R})\) of continuous numerical functions on \(\mathbf{X}\) ).
\item
  For every \(x \in \mathbf{X}\) , let \(\mathbf{f}(x)\) be the element of the dual \(\mathrm{F}'\) of \(\mathrm{F}\) defined by \(\langle z, \mathbf{f}(x) \rangle = z(x)\) for all \(z \in \mathrm{F}\) . Show that when \(\mathrm{F}'\) is equipped with the weak topology \(\sigma(\mathrm{F}', \mathrm{F})\) , \(\mathbf{f}\) is not \(\mu\) -measurable but \(\langle a, \mathbf{f} \rangle\) is \(\mu\) -measurable for every \(a \in \mathrm{F}\) .
\end{enumerate}

¶ 10) Let F be a Banach space, f a measurable mapping of X into F, such that the set A of \(x \in X\) where \(\mathbf{f}(x) \neq 0\) is a countable union of integrable sets. Show that there exists a sequence \((\mathbf{f}_n)\) of continuous functions with compact support, with values in F, such that the sequence \((\mathbf{f}_n(x))\) converges to \(\mathbf{f}(x)\) almost everywhere in X (note on the one hand that A is the union of a negligible set N and a sequence of pairwise disjoint compact sets \(K_n\) such that the restriction of f to each \(K_n\) is continuous; and on the other hand, that there exists a decreasing sequence \((\mathrm{U}_n)\) of open sets containing A such that \(\mu(\mathrm{U}_n \cap \mathbb{C}\mathrm{A})\) tends to 0 as n tends to infinity).

\begin{enumerate}
\def\labelenumi{\arabic{enumi})}
\setcounter{enumi}{10}
\tightlist
\item
  Let \(\mathbf{f}\) be a measurable mapping of \(X\) into a Banach space \(F\) . For every rational integer \(n\) (positive or negative) let \(A_{n}\) be the set of \(x \in X\) such that \(2^{n} \geqslant |\mathbf{f}(x)| > 2^{n-1}\) . For \(\mathbf{f}\) to be integrable, it is necessary and sufficient that the series with general term \(2^{n}\mu(A_{n})\) ( \(n \in \mathbf{Z}\) ) be convergent.
\end{enumerate}

¶ 12) Let \(\mu\) be a measure \(\geqslant 0\) on X. Let \((\mathbf{f}_n)\) be a sequence of integrable functions on X that converges pointwise on X to a function \(\mathbf{f}\) .

\begin{enumerate}
\def\labelenumi{\alph{enumi})}
\tightlist
\item
  Show that if \(\mathbf{f}\) is integrable and if
\end{enumerate}

\[
\int \mathbf {f} d \mu = \lim _ {n \rightarrow \infty} \int \mathbf {f} _ {n} d \mu ,
\]

then for every \(\varepsilon > 0\) there exist an integrable set \(A\) , an integrable function \(g \geqslant 0\) , and an integer \(n_0\) such that, for every \(n \geqslant n_0\) ,

\[
\left| \int \mathbf {f} _ {n} \varphi_ {\mathbb {C} A} d \mu \right| \leqslant \varepsilon
\]

and \(|\mathbf{f}_n(x)| \leqslant g(x)\) for all \(x \in A\) (consider an integrable set \(B\) such that

\[
\int | \mathbf {f} | \varphi_ {\mathbb {C} _ {\mathrm{B}}} d \mu \leqslant \varepsilon / 2
\]

and such that \(\mathbf{f}\) is bounded on B, and apply Egoroff's theorem).

\begin{enumerate}
\def\labelenumi{\alph{enumi})}
\setcounter{enumi}{1}
\item
  Suppose that, for every \(\varepsilon > 0\) , there exist a measurable set \(A\) , an integrable function \(g \geqslant 0\) and an integer \(n_0\) such that, for every \(n \geqslant n_0\) , one has \(\int |\mathbf{f}_n| \varphi_{\mathbb{C}A} d\mu \leqslant \varepsilon\) and \(|\mathbf{f}_n(x)| \leqslant g(x)\) for all \(x \in A\) . Show that, under these conditions, \(\mathbf{f}\) is integrable and \(\widetilde{\mathbf{f}}_n\) tends to \(\widetilde{\mathbf{f}}\) in the space \(L^1\) . Converse.
\item
  Suppose that \(\mathbf{F} = \mathbf{R}\) ; show by means of examples that the conditions of \(a\) ) are not sufficient, and that the conditions of \(b\) ) are not necessary, for \(\mathbf{f}\) to be integrable and for \(\int \mathbf{f} d\mu = \lim_{n \to \infty} \int \mathbf{f}_n d\mu\) to hold.
\end{enumerate}

¶ 13) Let \(\mu\) be a positive measure on X. If A is measurable, show that for every subset B of X,

\[
\mu^ {*} (\mathrm{B}) = \mu^ {*} (\mathrm{B} \cap \mathrm{A}) + \mu^ {*} (\mathrm{B} \cap \mathbb {C} \mathrm{A})
\]

(if \(\mu^{*}(\mathrm{B}) < +\infty\) , consider an integrable set \(\mathrm{B}_{1}\) such that \(\mathrm{B} \subset \mathrm{B}_{1}\) and \(\mu^{*}(\mathrm{B}) = \mu(\mathrm{B}_{1})\) (§4, Exer. 7 b))). Conversely, show that if A satisfies this condition, then A is measurable (cf.~§4, Exer. 6 e)).

\begin{enumerate}
\def\labelenumi{\arabic{enumi})}
\setcounter{enumi}{13}
\tightlist
\item
  Let \((\mathbf{A}_n)\) be a sequence of subsets of \(X\) such that, for every index \(n\) , there exists a measurable set \(\mathbf{B}_n \supset \mathbf{A}_n\) , the \(\mathbf{B}_n\) being pairwise disjoint. Show that
\end{enumerate}

\[
\mu^ {*} \left(\bigcup_ {n} \mathrm{A} _ {n}\right) = \sum_ {n} \mu^ {*} (\mathrm{A} _ {n}) \quad \text { and } \quad \mu_ {*} \left(\bigcup_ {n} \mathrm{A} _ {n}\right) = \sum_ {n} \mu_ {*} (\mathrm{A} _ {n}).
\]

\begin{enumerate}
\def\labelenumi{\arabic{enumi})}
\setcounter{enumi}{14}
\tightlist
\item
  Let \(\mu\) be a positive measure on \(X\) . A numerical function \(f\) (finite or not) defined on \(X\) is said to be quasi-integrable if it is measurable and if \(\mu^{*}(f) = \mu_{*}(f)\) (§4, Exer. 5). One then sets
\end{enumerate}

\[
\mu (f) = \mu^ {*} (f) = \mu_ {*} (f);
\]

again, one writes \(\int f d\mu\) in place of \(\mu(f)\) .

\begin{enumerate}
\def\labelenumi{\alph{enumi})}
\item
  Show that if \(f\) and \(g\) are quasi-integrable and if the sum \(\mu(f) + \mu(g)\) is defined, then \(f + g\) is defined almost everywhere and is quasi-integrable, and \(\mu(f + g) = \mu(f) + \mu(g)\) .
\item
  Deduce from a) that, for \(f\) to be quasi-integrable, it is necessary and sufficient that \(f\) be measurable and that at least one of the numbers \(\mu^{*}(f^{+}), \mu^{*}(f^{-})\) be finite; then, \(\mu(f) = \mu^{*}(f^{+}) - \mu^{*}(f^{-})\) .
\end{enumerate}

¶ 16) Let X be a compact space, \(\mu\) a positive measure on X. A bounded numerical function \(f\) defined on X is said to be continuous almost everywhere (for the measure \(\mu\) ) in X if the set of points of X where \(f\) is continuous (relative to X) has complement of measure zero.

\begin{enumerate}
\def\labelenumi{\alph{enumi})}
\item
  Give an example of a function \(f\) that is continuous almost everywhere and is such that there does not exist any continuous function \(g\) equal almost everywhere to \(f\) .
\item
  Suppose that the support of \(\mu\) is identical to X. Show that, for a bounded numerical function f defined on X to be equal almost everywhere to a function that is continuous almost everywhere in X, it is necessary and sufficient that there exist a subset H of X, whose complement is negligible, such that the restriction \(f|H\) of f to H is continuous (to prove that the condition is sufficient, note that H is dense in X, and that the extension of \(f|H\) to X that is lower semi-continuous on X (GT, IV, §6, Prop. 4) is a continuous function (relative to X) at every point of H). Deduce from this that f is measurable.
\item
  Deduce from b) that if in addition X is metrizable then, for every bounded numerical function f that is equal almost everywhere to a function that is continuous almost everywhere in X, there exists a sequence \((f_{n})\) of functions continuous on X, that is convergent at every point of X, and whose limit is almost everywhere equal to f (cf.~§4, Exer. 4 c)) (note that the proposition is true for a lower semi-continuous function f).
\item
  Let A be a perfect set without interior point, contained in X, and of measure \textgreater0 (cf.~§4, Exer. 4 a)). Show that there does not exist any function that is continuous almost everywhere in X and is equal almost everywhere to the upper semi-continuous function \(\varphi_{\mathrm{A}}\) .
\end{enumerate}

¶ 17) Let X be a compact space, \(\mu\) a positive measure on X. A set \(\mathcal{P}\) of finite partitions of X into integrable sets, directed for the relation \(\langle\varpi\) is coarser than \(\varpi'\rangle\) , is said to be fundamental if, for every entourage V of the uniform structure of X, there exists a partition \(\varpi = (A_i)\) of X belonging to \(\mathcal{P}\) such that all of the \(A_i\) are small of order V. For every finite partition \(\varpi = (A_k)\) belonging to \(\mathcal{P}\) , and every bounded numerical function \(f\) on X, set

\[
s _ {\varpi} (f) = \sum_ {k} \inf _ {x \in \mathrm{A} _ {k}} f (x) \cdot \mu (\mathrm{A} _ {k}) \quad \text { and } \quad S _ {\varpi} (f) = \sum_ {k} \sup _ {x \in \mathrm{A} _ {k}} f (x) \cdot \mu (\mathrm{A} _ {k})
\]

(the `Riemann sums' relative to \(f\) and the partition \(\varpi\) ).

\begin{enumerate}
\def\labelenumi{\alph{enumi})}
\item
  Show that \(s_{\varpi}(f) \leqslant \mu_{*}(f) \leqslant \mu^{*}(f) \leqslant S_{\varpi}(f)\) for every partition \(\varpi \in \mathcal{P}\) , and that \(s_{\varpi}(f)\) and \(S_{\varpi}(f)\) each tend to a limit with respect to the directed ordered set \(\mathcal{P}\) .
\item
  If \(\mathcal{P}\) is the (fundamental) set of all finite partitions of X into integrable sets, show that for every function \(f\) that is bounded and integrable, \(s_{\varpi}(f)\) and \(S_{\varpi}(f)\) tend to \(\int f d\mu\) with respect to \(\mathcal{P}\) .
\item
  If \(f\) is a bounded function that is continuous almost everywhere in \(X\) , then \(s_{\varpi}(f)\) and \(S_{\varpi}(f)\) tend to \(\int f d\mu\) with respect to every fundamental set \(\mathcal{P}\) of finite partitions of \(X\) into integrable sets (for every \(\varepsilon > 0\) , consider the closed set \(A\) of points where the oscillation of \(f\) is \(\geqslant \varepsilon\) , and for every partition \(\varpi \in \mathcal{P}\) whose sets are small of order \(V\) , consider separately those sets of \(\varpi\) that intersect \(V(A)\) and those that do not intersect it). If \(f\) is bounded and lower semi-continuous on \(X\) , show that \(s_{\varpi}(f)\) tends to \(\int f d\mu\) with respect to every fundamental set \(\mathcal{P}\) of finite partitions of \(X\) into integrable sets (regard \(f\) as the upper envelope of continuous functions).
\item
  A set \(A \subset X\) is said to be quadrable (for \(\mu\) ) if its characteristic function is continuous almost everywhere, or, what amounts to the same, if its boundary is \(\mu\) -negligible. Show that every point \(x_{0}\) of X has a fundamental system of quadrable open neighborhoods (for every neighborhood V of \(x_{0}\) , let f be a continuous function with values in {[}0,1{]}, equal to 1 at the point \(x_{0}\) and to 0 on CV; consider the sets of the x such that \(f(x) > \alpha\) , for \(0 < \alpha < 1\) ). From this, deduce that there exists a fundamental set P of finite partitions of X, such that every partition \(\varpi \in P\) consists of open sets and negligible sets.
\item
  Let \(\mathcal{P}\) be a fundamental set of finite partitions of \(X\) into integrable sets, such that every partition \(\varpi \in \mathcal{P}\) consists of open sets and negligible sets. For every bounded function \(f\) on \(X\) , let \(g\) be the largest lower semi-continuous function on \(X\) that is \(\leqslant f\) (GT, IV, §6, Prop. 4); show that for every partition \(\varpi \in \mathcal{P}\) , \(s_{\varpi}(f) = s_{\varpi}(g)\) . From this, deduce that for \(s_{\varpi}(f)\) and \(S_{\varpi}(f)\) to tend to a common limit with respect to \(\mathcal{P}\) , it is necessary that \(f\) be continuous almost everywhere in \(X\) .
\item
  Deduce from e) an example of a negligible function f and a fundamental set P of finite partitions of X into integrable sets, such that \(s_{\varpi}(f)\) and \(S_{\varpi}(f)\) do not tend to the same limit with respect to P (take for X the interval {[}0,1{]} of R, and for \(\mu\) the Lebesgue measure).
\end{enumerate}

¶ 18) Let X be a locally compact space such that, in X, the closure of every relatively compact open set is again an open set (a Stone space; cf.~Ch. II, §1, Exer. 13); let \(\mu\) be a positive measure on X with support identical to X and such that every bounded integrable numerical function on X is equivalent to a finite continuous function (§4, Exer. 10).

\begin{enumerate}
\def\labelenumi{\alph{enumi})}
\item
  Show that, in X, every nowhere dense set N is locally negligible (for every compact set K, consider the continuous function equivalent to \(\varphi_{K\cap \overline{N}}\) ).
\item
  Let \(f\) be a measurable numerical function (finite or not) on \(X\) , and let \(g\) be the largest lower semi-continuous function on \(X\) that is \(\leqslant f\) . Show that \(f\) and \(g\) are equal locally almost everywhere (note that if the restriction of \(f\) to a compact set \(K\) is continuous then \(f\) and \(g\) are equal on the interior of \(K\) , and make use of \(a\) ).
\item
  Deduce from b) that in X, every set locally negligible for \(\mu\) is a nowhere dense set. In particular, every meager set in X is nowhere dense.
\end{enumerate}

\begin{enumerate}
\def\labelenumi{\arabic{enumi})}
\setcounter{enumi}{18}
\item
  Let \(\mu\) be a positive measure on a locally compact space X. Let f be a \(\mu\) -measurable mapping of X into a complete metric space F. In order that, in a compact subset K of X, f may be uniformly approximated by measurable step functions, it is necessary and sufficient that \(f(\mathrm{K})\) be relatively compact in F.
\item
  Let \(X\) be a compact space, \(\mu\) a positive measure on \(X\) , and \((f_n)\) a sequence of \(\mu\) -measurable numerical functions. Show that the following properties are equivalent:
\end{enumerate}

\begin{enumerate}
\def\labelenumi{(\roman{enumi})}
\item
  there exists a subsequence \((f_{n_k})\) of \((f_n)\) tending to 0 almost everywhere in \(X\) ;
\item
  there exists a sequence \((\lambda_{n})\) of finite real numbers such that
\end{enumerate}

\[
\limsup _ {n \to \infty} | \lambda_ {n} | > 0
\]

and such that the series with general term \(\lambda_{n}f_{n}(x)\) converges almost everywhere in X; (iii) there exists a sequence \((\lambda_n)\) of finite real numbers such that

\[
\sum_ {n = 1} ^ {\infty} | \lambda_ {n} | = + \infty
\]

and such that the series with general term \(\lambda_{n}f_{n}(x)\) is absolutely convergent almost everywhere in X.

(To see that (i) implies (ii) and (iii), make use of Egoroff's theorem. To see that (iii) implies (i), show that (iii) implies the existence of an increasing sequence \((\mathbf{A}_k)\) of measurable subsets of X and a subsequence \((f_{n_k})\) of \((f_n)\) , such that \(\mu (\mathbf{A}_k)\) tends to \(\mu (\mathbf{X})\) and \(\int |f_{n_k}\varphi_{\mathrm{A}_k}|d\mu\) tends to 0.)

\begin{enumerate}
\def\labelenumi{\arabic{enumi})}
\setcounter{enumi}{20}
\item
  Let X be a locally compact space, \(\mu\) a positive measure on X, and \((\mathrm{U}_{\alpha})_{\alpha\in\mathbf{A}}\) an open covering of X. For every \(\alpha\in A\) , let \(f_{\alpha}\) be a mapping of \(U_{\alpha}\) into a set G. Assume that for every pair of indices \(\alpha,\beta\) , the set of points \(x\in U_{\alpha}\cap U_{\beta}\) such that \(f_{\alpha}(x)\neq f_{\beta}(x)\) is locally \(\mu\) -negligible. Show that there exists a mapping f of X into G such that, for every \(\alpha\in A\) , the set of \(x\in U_{\alpha}\) such that \(f(x)\neq f_{\alpha}(x)\) is locally \(\mu\) -negligible. (Consider first the case that X is compact, by covering X by a finite number of sets \(\mathrm{U}_{\alpha}\) ; pass to the general case with the help of Prop. 14 of No.~9.)
\item
  Show that if the positive measure \(\mu\) is such that there exist open sets of arbitrarily small measure \(>0\) then, for every \(a > 0\) , the topology induced on the set of \(\mathbf{f} \in \mathcal{L}_{\mathrm{F}}^{p}\) such that \(\mathrm{N}_p(\mathbf{f}) \leqslant a\) by the topology of convergence in measure is strictly coarser than the topology of convergence in mean of order \(p\) .
\item
  Let \((\mathbf{f}_{n})\) be a sequence of functions in \(S_{F}\) such that, for every integrable set A and every subsequence \((\mathbf{f}_{n_{k}})\) of \((\mathbf{f}_{n})\) , there exists a subsequence of \((\mathbf{f}_{n_{k}})\) that converges to 0 almost everywhere in A; show that the sequence \((\mathbf{f}_{n})\) converges in measure to 0. (Argue by contradiction.)
\end{enumerate}

¶ 24) Let X be the interval {[}0,1{]} of R, μ the Lebesgue measure on X. Show that for every Banach space F, every continuous linear form on \(S_F\) is identically zero. (For every neighborhood V of 0 in \(S_F\) , show that there exists an integer n \textgreater{} 0 such that V + V + \ldots{} + V (n terms) contains a line.)

\begin{enumerate}
\def\labelenumi{\arabic{enumi})}
\setcounter{enumi}{24}
\tightlist
\item
  \begin{enumerate}
  \def\labelenumii{\alph{enumii})}
  \tightlist
  \item
    For a subset H of \(S_{F}\) to be precompact, it is necessary and sufficient that, for every \(\varepsilon > 0\) and every integrable set \(A \subset X\) , there exist a compact set \(M \subset F\) and a partition of A into a finite number of integrable sets \(A_{i}\) such that, for every \(f \in H\) , there exists an integrable set \(B \subset A\) of measure \(\mu(B) \leqslant \varepsilon\) and having the following properties: \(1^{\circ}\) every point of \(f(A - B)\) has distance \(\leqslant \varepsilon\) from M; \(2^{\circ}\) in each of the sets \(A_{i} \cap C B\) , the oscillation of f is \(\leqslant \varepsilon\) .
  \end{enumerate}
\end{enumerate}

\begin{enumerate}
\def\labelenumi{\alph{enumi})}
\setcounter{enumi}{1}
\item
  Show that for a subset H of \(\mathcal{L}_{\mathrm{F}}^{1}\) to be relatively quasi-compact, it is necessary and sufficient that it be precompact in \(\mathcal{S}_{\mathrm{F}}\) and equi-integrable.
\item
  For a sequence \((\mathbf{f}_{n})\) of functions in \(L_{F}^{1}\) to be a Cauchy sequence for the topology of convergence in mean, show that it is necessary and sufficient that \((\mathbf{f}_{n})\) be a Cauchy sequence for the topology of convergence in measure and that the set of the \(f_{n}\) be equi-integrable.
\end{enumerate}

\(d\) ) Extend these properties to the spaces \(\mathcal{L}_{\mathrm{F}}^{p}\) for \(p > 1\)

\begin{enumerate}
\def\labelenumi{\arabic{enumi})}
\setcounter{enumi}{25}
\tightlist
\item
  Let \(f\) be a numerical function defined on \(\mathbf{R}\) . Show that if, for every \(x \in \mathbf{R}\) , \(\liminf_{y \to x, y \geqslant x} f(y) \geqslant f(x)\) , then \(f\) is \(\mu\) -measurable for every measure \(\mu\) on \(\mathbf{R}\) .
\end{enumerate}

¶ 27) Let X be a locally compact space, \(\mu\) a real measure on X, K a compact subset of X, \((f_n)\) a sequence of bounded \(\mu\) -measurable numerical functions defined on K and separating the points of K. Show that for every bounded \(\mu\) -measurable numerical function g defined on K, and every \(\varepsilon > 0\) , there exist a polynomial \(h = \mathrm{P}((f_n))\) in the \(f_n\) , and a compact subset \(\mathrm{K}_1\) of K, such that \(\|h\| \leqslant 2\|g\|\) , \(|\mu|(\mathrm{K} - \mathrm{K}_1) \leqslant \varepsilon\) , and, for every \(x \in \mathrm{K}_1\) , \(|h(x) - g(x)| \leqslant \varepsilon\) . (Consider the mapping \(x \mapsto (f_n(x))\) of K into \(\mathbf{R}^{\mathbf{N}}\) and the closure of its image, and apply the Stone--Weierstrass theorem suitably.) From this, deduce that for \(1 \leqslant p < +\infty\) , the set of polynomials in the \(f_n\) is dense in \(\mathcal{L}^p(\mathrm{K})\) . Are these properties again true when the sequence \((f_n)\) is replaced by an uncountable family of bounded measurable functions separating the points of K?

\begin{enumerate}
\def\labelenumi{\arabic{enumi})}
\setcounter{enumi}{27}
\tightlist
\item
  Let X be a locally compact space, \(\mu\) a positive measure on X, and P the set of numerical functions (finite or not) \(f \geqslant 0\) defined on X and \(\mu\) -measurable. Let \(\lambda\) be a function defined on P with values \(\geqslant 0\) (finite or not) that is positively homogeneous, increasing and convex (cf.~Ch. I, No.~1) and, moreover, such that: \(1^{\circ} \lambda(f) = 0\) for every \(\mu\) -negligible function \(f \geqslant 0\) ; \(2^{\circ}\) for every increasing sequence \((f_{n})\) of functions in P, \(\lambda(\sup f_{n}) = \sup_{n} \lambda(f_{n})\) . For every Banach space F, denote by \(L_{F}^{\lambda}\) the set of \(\mu\) -measurable mappings f of X into F such that \(\lambda(|\mathbf{f}|)\) is finite. Show that \(\lambda(|\mathbf{f}|)\) is a semi-norm on \(L_{F}^{\lambda}\) , and that \(L_{F}^{\lambda}\) is complete for the topology defined by this semi-norm.
\end{enumerate}

¶ 29) Let \(\mu\) be a measure on a locally compact space X. For every \(\mu\) -measurable numerical function \(f \geqslant 0\) defined on X, the mapping \(f_{\mu}' : t \mapsto |\mu|^*(f([t, +\infty]))\) of \(\mathbf{R}_+\) into \(\overline{\mathbf{R}}_{+}\) is decreasing and right-continuous. If \(\nu\) is a measure on a second locally compact space Y and \(g\) is a \(\nu\) -measurable numerical function \(\geqslant 0\) defined on Y, \(f\) and \(g\) are said to be equimeasurable (for \(\mu\) and \(\nu\) respectively) if \(f_{\mu}^{\prime} = g_{\nu}^{\prime}\) . One denotes by \(f^{*}\) (or \(f_{\mu}^{*}\) ) the function defined on \(\mathbf{R}_{+}\) that is equal, for all \(s \in \mathbf{R}_{+}\) , to the supremum in \(\overline{\mathbf{R}}_{+}\) of the set of numbers \(a\) such that \(f_{\mu}^{\prime}(a) \geqslant s\) (the supremum being equal to 0 if this set is empty).

\begin{enumerate}
\def\labelenumi{\alph{enumi})}
\item
  Show that the function \(f^{*}\) is decreasing, left-continuous in \(\mathbf{R}_{+}\) , and equi-measurable to \(f\) (for the measure \(\mu\) and Lebesgue measure on \(\mathbf{R}_{+}\) ); \(f^{*}\) is called the decreasing rearrangement of \(f\) .
\item
  If \(0 \leqslant f \leqslant g\) are two \(\mu\) -measurable functions on \(X\) , show that \(f^* \leqslant g^*\) . If \((f_n)\) is an increasing sequence of \(\mu\) -measurable functions \(\geqslant 0\) on \(X\) , and \(f = \sup_{n} f_n\) , show
\end{enumerate}

that \(f^{*}(s) = \sup_{n}f_{n}^{*}(s)\) at all the points where \(f^{*}\) is continuous.

\begin{enumerate}
\def\labelenumi{\alph{enumi})}
\setcounter{enumi}{2}
\item
  Show that for every \(\mu\) -measurable numerical function \(f \geqslant 0\) , \(|\mu|^{*}(f) = \int^{*} f^{*}(s) ds\) . (Consider first the case of a step function, then make use of \(b\) ).)
\item
  Let w be a numerical function defined on \(R_{+}\) that is finite (except possibly at the point s = 0) and decreasing. For every \(\mu\) -measurable numerical function \(f \geqslant 0\) defined on X, set
\end{enumerate}

\[
\lambda (f) = \left(\int^ {*} w (s) \bigl (f ^ {*} (s) \bigr) ^ {p} d s\right) ^ {1 / p} \quad (1 \leqslant p <   + \infty).
\]

Show that this function satisfies the conditions of Exer. 28. (To prove that it is convex, consider first the case that w is a decreasing step function, using c), then pass to the limit to treat the general case.)

\begin{enumerate}
\def\labelenumi{\alph{enumi})}
\setcounter{enumi}{4}
\tightlist
\item
  With the notations of \(d\) ), write \(\mathcal{L}_{\mathrm{F}}^{p,w}\) instead of \(\mathcal{L}_{\mathrm{F}}^{\lambda}\) , and set \(\mathrm{N}_{p,w}(\mathbf{f}) = \lambda (|\mathbf{f}|)\) for every \(\mu\) -measurable mapping \(\mathbf{f}\) of X into F. Show that if \(\mathbf{f} \in \mathcal{L}_{\mathrm{F}}^{p,w}\) and if, for every \(n\) , \(\mathbf{f}_n\) denotes the function equal to \(\mathbf{f}\) if \(|\mathbf{f}| \leqslant n\) and to \(nf / |\mathbf{f}|\) if \(|\mathbf{f}| > n\) , then the sequence \((|\mathbf{f}_n|^*)\) tends almost everywhere to \(|\mathbf{f}|^*\) and one has \(\lim_{n \to \infty} \mathrm{N}_{p,w}(\mathbf{f} - \mathbf{f}_n) = 0\) . If \(\mathbf{f}\) is \(\mu\) -negligible, then \(\mathrm{N}_{p,w}(\mathbf{f}) = 0\) .
\end{enumerate}

¶ 30) Let \(\Phi(t, u, v)\) be a finite and continuous numerical function for \(0 \leqslant t \leqslant 1\) , \(u \geqslant 0\) , \(v \geqslant 0\) . For every numerical function \(f \geqslant 0\) defined on I = {[}0,1{]}, measurable for Lebesgue measure, \(f^*\) denotes the decreasing rearrangement of \(f\) (Exer. 29). In order that, for every pair of functions \(f \geqslant 0\) , \(g \geqslant 0\) , measurable (for Lebesgue measure) and bounded on I, one have

\[
\int_ {0} ^ {1} \Phi (t, f (t), g (t)) d t \leqslant \int_ {0} ^ {1} \Phi (t, f ^ {*} (t), g ^ {*} (t)) d t,\tag{1}
\]

it is necessary and sufficient that the function \(\Phi\) satisfy the following three conditions:

\[
\Phi (t, u + h, v + h) - \Phi (t, u + h, v) - \Phi (t, u, v + h) + \Phi (t, u, v) \geqslant 0\tag{2}
\]

for all \(t \in I, u \geqslant 0, v \geqslant 0\) and \(h \geqslant 0\) ;

\[
\begin{array}{r l} \int_ {0} ^ {\delta} & (\Phi (a + \delta + t, u, v) - \Phi (a + \delta + t, u + h, v) \\ & + \Phi (a + t, u + h, v) - \Phi (a + t, u, v)) d t \geqslant 0 \end{array}\tag{3}
\]

and

\[
\begin{array}{l} \int_ {0} ^ {\delta} \big (\Phi (a + \delta + t, u, v) - \Phi (a + \delta + t, u, v + h) \\ \qquad \qquad \qquad + \Phi (a + t, u, v + h) - \Phi (a + t, u, v) \big) d t \geqslant 0 \end{array}\tag{4}
\]

for all \(a, h, \delta\) such that \(0 \leqslant a \leqslant 1 - 2\delta\) , \(h \geqslant 0\) .

(To prove that (3) and (4) are necessary, take \(f\) and \(g\) to be suitable step functions; next deduce (2) from (3) and (4). To prove that the conditions are sufficient, first deduce from (2), using the continuity of \(\Phi\) , that

\[
\Phi (t, u + h, v + k) - \Phi (t, u + h, v) - \Phi (t, u, v + k) + \Phi (t, u, v) \geqslant 0\tag{5}
\]

for all \(t \in \mathrm{I}\) , \(u \geqslant 0\) , \(v \geqslant 0\) , \(h \geqslant 0\) , \(k \geqslant 0\) . Then using (3) and (4), prove (1) when \(f\) and \(g\) are step functions whose points of discontinuity in \(\mathrm{I}\) are of the form \(r / n\) ( \(0 \leqslant r \leqslant n\) ), by comparing successively the integrals

\[
\int_ {0} ^ {1} \Phi (t, f _ {i} (t), g _ {i} (t)) d t
\]

and

\[
\int_ {0} ^ {1} \Phi (t, f _ {j} (t), g _ {j} (t)) d t,
\]

where \(f_{i}\) and \(g_{i}\) differ from \(f_{j}\) and \(g_{j}\) only in two consecutive intervals \([(r - 1) / n, r / n]\) and \([r / n, (r + 1) / n]\) . Finally, pass to the limit in a suitable way to prove (1) in the general case.)

If \(\Phi\) is twice continuously differentiable, the conditions (2), (3) and (4) are equivalent, respectively, to

\[
\frac {\partial^ {2} \Phi}{\partial u \partial v} \geqslant 0, \quad \frac {\partial^ {2} \Phi}{\partial t \partial u} \leqslant 0, \quad \frac {\partial^ {2} \Phi}{\partial t \partial v} \leqslant 0.
\]

Generalize to functions \(\Phi(t, u_1, \ldots, u_m)\) of any number of variables.

\setcounter{exercisegroup}{6}
\edef\CurrentExerciseGroup{\arabic{exercisegroup}}
